\documentclass[10pt,aspectratio=169]{beamer}
%==============================================================================
%  Running MfMax v1 -- slide deck
%  Companion to doc/mfmax_v1_howto.tex (the full how-to).
%  Source of record: indirect/mfmax/MFMAX_V1_RUNBOOK.md
%  All numbers measured on this machine 2026-08-15.
%==============================================================================
\usetheme{default}
\usecolortheme{dolphin}
\setbeamertemplate{navigation symbols}{}
\setbeamertemplate{footline}[frame number]
\setbeamercolor{frametitle}{fg=black}
\setbeamerfont{frametitle}{series=\bfseries}
\setbeamertemplate{itemize item}{\textendash}
\setbeamertemplate{itemize subitem}{\textperiodcentered}
\setbeamercolor{block title}{bg=gray!18,fg=black}
\setbeamercolor{block body}{bg=gray!8}
\setbeamercolor{block title alerted}{bg=orange!25,fg=black}
\setbeamercolor{block body alerted}{bg=orange!10}

\usepackage{amsmath,amssymb}
\usepackage{booktabs}
\usepackage{xcolor}
\usepackage{listings}

\definecolor{codegray}{gray}{0.94}
\definecolor{codekw}{rgb}{0.13,0.13,0.7}
\definecolor{codecomment}{rgb}{0.0,0.5,0.0}
\definecolor{codestr}{rgb}{0.6,0.1,0.1}

\lstdefinestyle{sh}{
  language=bash, basicstyle=\ttfamily\scriptsize,
  keywordstyle=\color{codekw}, commentstyle=\color{codecomment}\itshape,
  stringstyle=\color{codestr}, backgroundcolor=\color{codegray},
  numbers=none, breaklines=true, frame=none,
  showstringspaces=false, columns=fullflexible, keepspaces=true,
  aboveskip=4pt, belowskip=4pt}
\lstdefinestyle{plain}{
  basicstyle=\ttfamily\tiny, backgroundcolor=\color{codegray},
  numbers=none, breaklines=true, frame=none,
  showstringspaces=false, columns=fullflexible, keepspaces=true,
  aboveskip=4pt, belowskip=4pt}
\lstdefinestyle{py}{
  language=Python, basicstyle=\ttfamily\scriptsize,
  keywordstyle=\color{codekw}, commentstyle=\color{codecomment}\itshape,
  stringstyle=\color{codestr}, backgroundcolor=\color{codegray},
  numbers=none, breaklines=true, frame=none,
  showstringspaces=false, columns=fullflexible, keepspaces=true,
  aboveskip=4pt, belowskip=4pt}
\lstset{style=sh}

\newcommand{\code}[1]{\texttt{\small #1}}
\newcommand{\file}[1]{\texttt{\small #1}}
\newcommand{\Tmax}{T_{\max}}
\newcommand{\Lf}{L_f}

\title{\bfseries Running MfMax v1}
\subtitle{Build, drive and read the Gergaud-group indirect minimum-fuel solver}
\author{\file{earth\_elliptic\_to\_geo/indirect/mfmax}}
\date{2026-08-23}

\begin{document}

%------------------------------------------------------------------
\begin{frame}
\titlepage
\begin{center}\small
Full how-to: \file{doc/mfmax\_v1\_howto.tex} \quad$\cdot$\quad
Source of record: \file{indirect/mfmax/MFMAX\_V1\_RUNBOOK.md}\\[2pt]
All numbers measured on this machine, 2026-08-15
\end{center}
\end{frame}

%------------------------------------------------------------------
\begin{frame}{What MfMax is, and why we keep it}

\begin{columns}[T]
\begin{column}{0.52\textwidth}
\textbf{What}
\begin{itemize}
\item Gergaud group's Fortran solver, ENSEEIHT-IRIT 2004
\item Low-thrust \textbf{minimum-fuel} orbit transfer in modified equinoctial
      elements
\item \textbf{Indirect}: PMP single shooting $+$ HOMPACK90 differential
      homotopy
\item Solves \emph{exactly} the transfer our
      \file{earth\_elliptic\_to\_geo} campaign reproduces by direct
      collocation
\end{itemize}
\end{column}
\begin{column}{0.44\textwidth}
\textbf{Why we care}
\begin{itemize}
\item Our one genuinely \textbf{independent indirect cross-check}
\item Built + validated here: \textbf{0.7 s} to a converged transfer
\item Its ideas are borrowable (last slide)
\item Its converged costates are clean shooting seeds
\end{itemize}
\end{column}
\end{columns}

\vspace{4pt}
\begin{block}{The whole talk in one line}
\code{make} with two legacy flags, drop an \file{in.dat} in a directory, run
\code{path}, then check four things that the program does \emph{not} check for
you.
\end{block}
\end{frame}

%------------------------------------------------------------------
\begin{frame}{The problem it solves}

With MEE state $x = (P,e_x,e_y,h_x,h_y)$, mass $m$, control $u$ in the unit
ball ($T = \Tmax u$):
\begin{equation*}
\min_{\|u\|\le 1}\ \int_{0}^{t_f}\!\|u\|\,\mathrm{d}t
\qquad\text{s.t.}\qquad
\dot x = f(x,u),\quad \dot m = -\beta\,\Tmax\|u\|
\end{equation*}

\vspace{2pt}
Benchmark case (HMG-2004), all units \textbf{Mm, hours, kg}:

\begin{center}\small
\begin{tabular}{@{}ll@{}}
\toprule
depart & $P = 11.625$ Mm, $e = 0.75$, $i = 7^\circ$, $m_0 = 1500$ kg,
$L_0 = \pi$\\
arrive & GEO: $P = 42.165$ Mm, circular, equatorial\\
engine & $\Tmax = 10$ N, $\mathrm{Isp} = 1994.75$ s\\
\bottomrule
\end{tabular}
\end{center}

\begin{alertblock}{Engine constant --- before you compare masses}
\code{par(3)}$= 0.0142$ h/Mm $\Rightarrow v_e = 19\,561.8$ m/s
$\Rightarrow \mathrm{Isp} = 1994.75$ s. That is the benchmark's \emph{exact}
value; our direct campaign defaults to a round $2000$ s. A $0.27\%$
difference.
\end{alertblock}
\end{frame}

%------------------------------------------------------------------
\begin{frame}{The v1 device: fixed longitude span, free final time}

\begin{itemize}
\item Independent variable is the \textbf{longitude} $L$, integrated
      $L_0 \to \Lf$ \hfill (Sundman-like)
\item $\Lf$ is \textbf{fixed}: $\;\Lf = L_0 + \mathrm{par}(1)\,(L_{\min}-L_0)$
\item $t_f$ is \textbf{free} --- carried as state $x_8$ with
      $\mathrm{d}x_8/\mathrm{d}L \equiv 0$
\item Normalized time $s = t/t_f$ is state $x_6$; the condition
      $s(\Lf) = 1$ closes the system
\end{itemize}

\begin{center}\small
\begin{tabular}{@{}llllllllc@{}}
\toprule
index & 1 & 2 & 3 & 4 & 5 & 6 & 7 & 8 \quad|\quad 9--16\\
\midrule
meaning & $P$ & $e_x$ & $e_y$ & $h_x$ & $h_y$ & $s = t/t_f$ & $m$ &
$t_f$ \quad|\quad costates\\
\bottomrule
\end{tabular}
\end{center}

\begin{block}{The transferable trick}
Free an unknown scalar by carrying it as a \textbf{zero-derivative state} and
closing with a \textbf{normalized terminal condition}. Worth stealing
independently of MfMax.
\end{block}

\begin{alertblock}{Slot 6 is aliased --- read \file{ftir.f90:74--102} first}
What \code{Ufun}/\code{Dhfun} see as $x_6$ is overwritten with the longitude
$l$; the \emph{integrated} slot 6 is $s$, passed separately as \code{t}.
\code{Phifun} then scales every derivative by
$\mathrm{d}s/\mathrm{d}L = 1/F(6)$.
\end{alertblock}
\end{frame}

%------------------------------------------------------------------
\begin{frame}{\code{lambda} is a pair --- two unrelated jobs}

\begin{center}\small
\begin{tabular}{@{}llll@{}}
\toprule
 & drives & $0$ means & $1$ means\\
\midrule
$\lambda_1$ & initial-condition homotopy & start \emph{at the target orbit} &
the true departure orbit\\
$\lambda_2$ & energy $\to$ fuel & smooth, energy-like control &
true bang-bang minimum fuel\\
\bottomrule
\end{tabular}
\end{center}

\vspace{2pt}
\textbf{$\lambda_1$} (\code{homCI}, discrete continuation):
$y(1{:}5) = (1-\lambda_1)\,\mathrm{par}(5{:}9) + \lambda_1\,y_0(1{:}5)$.
At $\lambda_1 = 0$ you start \emph{on the target orbit} --- a trivial problem.
Step \code{par(10)}, halved on failure down to \code{par(11)}.

\textbf{$\lambda_2$} (\code{Fullpath}, HOMPACK arclength):
\begin{equation*}
\min \int_0^{t_f}\Big[\lambda_2\|u\| + (1-\lambda_2)\|u\|^2\Big]\mathrm{d}t
\qquad
\rho =
\begin{cases}
0, & \|\Phi\| \le \sigma\\
\frac{\|\Phi\|-\sigma}{2(1-\lambda_2)}, & \text{ramp}\\
1, & \|\Phi\| \ge \sigma + 2(1-\lambda_2)
\end{cases}
\end{equation*}
with $\sigma = \lambda_2 - \beta\Tmax p_m$. At $\lambda_2 = 1$ the ramp
collapses: exactly bang-bang.

\begin{block}{}
Same energy-to-fuel idea as our $\varepsilon$-homotopy --- but by
\textbf{arclength continuation with automatic steplength}, not a
hand-scheduled ladder.
\end{block}
\end{frame}

%------------------------------------------------------------------
\begin{frame}[fragile]{Quick start --- five commands, 90 seconds}

\begin{lstlisting}
MFMAX=~/Desktop/optimal_control/orbit_transfer/earth_elliptic_to_geo/indirect/mfmax

export SDKROOT=$(xcrun --show-sdk-path)
cp -R $MFMAX/mfmax-v1 /tmp/mfmax-v1
cd /tmp/mfmax-v1 && rm -f src/*.o src/*.mod
make -C src F90=gfortran OPTF="-O3 -c -std=legacy -fallow-argument-mismatch" \
     OPTO="-O3 -o"
mkdir -p run10N && cd run10N && cp ../matlab/IN/in10N.dat in.dat && ../bin/path
\end{lstlisting}

\begin{block}{Success looks like this}
\begin{lstlisting}[style=plain]
Final mass (kg) =    1381.4074215554372
Final time (h)  =    146.39467958103148
** Total execution time  (secs) =   0.69800000000395812
\end{lstlisting}
\end{block}

\textbf{Prerequisites:} gfortran (verified 13.1.0) and Xcode command line
tools. No MATLAB, no CasADi.
\end{frame}

%------------------------------------------------------------------
\begin{frame}[fragile]{Step 1--2 --- copy out of tree, then build}

\textbf{Build out of tree.} The repository copy has eight \file{.o} and four
\file{.mod} files from a 2007 build \emph{committed into it}, and \file{bin/}
is empty.

\begin{lstlisting}
cp -R $MFMAX/mfmax-v1 /tmp/mfmax-v1 && cd /tmp/mfmax-v1
rm -f src/*.o src/*.mod        # else make skips recompiles
export SDKROOT=$(xcrun --show-sdk-path)
make -C src F90=gfortran OPTF="-O3 -c -std=legacy -fallow-argument-mismatch" \
     OPTO="-O3 -o"
\end{lstlisting}

\begin{itemize}
\item \code{F90=gfortran} overrides the makefile's \code{g95}, which is long
      dead. \textbf{No source edit needed.}
\item \code{-std=legacy -fallow-argument-mismatch} are \textbf{required}:
      HOMPACK90 and the LAPACK/rkf45 slices are F77 with mismatched argument
      types that gfortran $\ge 10$ rejects.
\item Expect only noise warnings: clang deployment version, duplicate
      \code{-lemutls\_w}/\code{-lgcc}, ``built for newer macOS''.
\item Result: 8 objects $+$ \file{bin/path}.
\end{itemize}
\end{frame}

%------------------------------------------------------------------
\begin{frame}[fragile]{Step 3 --- run}

\code{path} takes \textbf{no arguments}. It reads \file{./in.dat} from the
current directory and writes beside it --- so give every case its own
directory.

\begin{lstlisting}
mkdir run10N && cd run10N
cp ../matlab/IN/in10N.dat in.dat
../bin/path | tee run.log
\end{lstlisting}

\begin{center}\small
\begin{tabular}{@{}ll@{}}
\toprule
\file{out.dat} & header $+$ $\mathrm{NIT}+1$ trajectory records\\
\file{next.dat} & restart file in \file{in.dat} format, $\lambda = (1,1)$\\
\file{fort.2} & HOMPACK arclength trace (unit $=$ the \code{TRACE} value)\\
\bottomrule
\end{tabular}
\end{center}

\begin{block}{Warm start}
\file{next.dat} \emph{is} an \file{in.dat}. \code{cp next.dat in.dat}, change
what you want, rerun. Strongly preferred for a nearby case.
\end{block}
\end{frame}

%------------------------------------------------------------------
\begin{frame}{Step 4 --- verify (the program will not do this for you)}

\begin{enumerate}
\item \textbf{Did it reach $\lambda = (1,1)$?} The \code{AFTER REFINEMENT}
      block must show both entries at $1.0$. Anything less and you solved a
      \emph{regularized} problem, not minimum fuel.
\item \textbf{Is the residual small?} \code{|| Ftir(Ysol) ||} at the
      $10^{-9}$ level. Reference run: $2.24\times10^{-9}$.
\item \textbf{Did it hit the target?} \emph{Not printed} --- read
      \file{out.dat}. Reference: $P = 42.164999999$ Mm, i.e.
      $8.9\times10^{-10}$ Mm error; $e = 1.4\times10^{-11}$;
      $h_x,h_y \sim 10^{-13}$.
\item \textbf{Do the invariants hold?} $\|u\| \in \{0,1\}$ exactly (genuine
      bang-bang, so $\lambda_2 = 1$ really was attained); $x_8$ constant
      (measured spread \emph{exactly} $0$); $s$ monotone $0\to1$; mass
      monotone decreasing.
\end{enumerate}

\begin{alertblock}{And check where the \code{IFAIL} lines sit --- see the
gotchas slide}
\end{alertblock}
\end{frame}

%------------------------------------------------------------------
\begin{frame}{The reference run --- reproduce this before trusting anything}

Stock \file{in10N.dat}, unmodified. $\Lf = \pi + 2.0(29.9012-\pi) = 56.6608$
rad fixed, $t_f$ free.

\begin{center}
\begin{tabular}{@{}ll@{}}
\toprule
wall time & $0.70$ s ($0.47$ s in the differential homotopy)\\
Jacobian evaluations & 18\\
shooting residual $\|F\|$ & $2.24\times10^{-9}$\\
\textbf{final mass} & \textbf{1381.4074 kg} \; (118.5926 kg propellant)\\
\textbf{final time} (the free unknown) & \textbf{146.3947 h $=$ 6.0998 d}\\
revolutions & 8.5178\\
thrust structure & 18 arcs / 17 switches, 19.0\% burn duty\\
\bottomrule
\end{tabular}
\end{center}

\begin{block}{Plausibility against our own campaign}
Our certified direct run gives 1377.10 kg at 7.326 rev ($c_{t_f}=1.5$); its
$c_{t_f}$ front runs 1377--1385 kg over $c_{t_f}\in[2.0,2.5]$. MfMax's
1381.41 kg at 8.518 rev sits inside that band, in the right direction --- a
longer transfer needs less propellant.
\end{block}
\end{frame}

%------------------------------------------------------------------
\begin{frame}[fragile]{\file{in.dat} --- twelve records}

\begin{columns}[T]
\begin{column}{0.46\textwidth}
\begin{lstlisting}[style=plain]
 10. 0. 0. 0. 0. 0. 0. 0.
 8   9  10  11  12  13  14 15
 11.625 0.75 0. 0.0612 0. 0. 1500. 0. ...
 3.14159
 29.901174314034
 0. 0.
 2. 10. 0.0142 0.516586E+04 42.165 ...
  1
 1000
 1.0E-4
 2
 0. 0.
\end{lstlisting}
\end{column}
\begin{column}{0.52\textwidth}
\scriptsize
1 \code{ZI(1:8)} initial guess for the unknowns\\
2 \code{FREE0(1:8)} which \code{Y0} slots they occupy\\
3 \code{Y0(1:16)} initial state $+$ costate\\
4 \code{L0} initial longitude [rad]\\
5 \code{Lmin} reference final longitude\\
6 \code{lambda(1:2)} homotopy start\\
7 \code{PAR(1:11)} (next slide)\\
8 \code{IPAR(1)} run mode\\
9 \code{NIT} trajectory samples\\
10 \code{jac\_step} finite-difference step\\
11 \code{TRACE} \textbf{a unit number}, not a verbosity\\
12 \code{SSPAR(1:2)} HOMPACK steplength (0 $=$ auto)
\end{column}
\end{columns}

\vspace{3pt}
\code{FREE0 = [8..15]}: the eight unknowns are $t_f$ plus the seven costates
$p_P,p_{e_x},p_{e_y},p_{h_x},p_{h_y},p_s,p_m$ at $L_0$. So \textbf{\code{ZI(1)}
is the $t_f$ guess in hours} --- forced to 10 if you pass $\le 0$.
\end{frame}

%------------------------------------------------------------------
\begin{frame}{\code{PAR} and \code{IPAR}}

\begin{center}\small
\begin{tabular}{@{}l p{9.4cm} c@{}}
\toprule
slot & meaning & 10 N\\
\midrule
\code{par(1)} & \textbf{longitude-span multiplier}
$\Lf = L_0 + \mathrm{par}(1)(L_{\min}-L_0)$ --- v1's analogue of $c_{t_f}$ &
2.0\\
\code{par(2)} & max thrust [N] (internally $\times 12.96$ to kg\,Mm/h$^2$) &
10.\\
\code{par(3)} & $\beta$ mass-flow coefficient [h/Mm],
$\dot m = -\beta\Tmax\|u\|$ & 0.0142\\
\code{par(4)} & $\mu$ [Mm$^3$/h$^2$] & 5165.86\\
\code{par(5:9)} & \textbf{target orbit} $(P,e_x,e_y,h_x,h_y)$ --- terminal
constraint \emph{and} the $\lambda_1{=}0$ easy start & 42.165 0 0 0 0\\
\code{par(10:11)} & initial / minimum $\lambda_1$ continuation step &
0.1, 0.01\\
\bottomrule
\end{tabular}
\end{center}

\begin{center}\small
\begin{tabular}{@{}cl@{}}
\toprule
\code{IPAR(1)} & behaviour\\
\midrule
$\le 1$ & full homotopy (\code{Fullpath}) --- \textbf{the normal mode}\\
$2$ & no solve; integrate from the given \code{ZI} and write the trajectory\\
$3$ & single shooting at the given \code{lambda}, print residual, write\\
\bottomrule
\end{tabular}
\end{center}
\end{frame}

%------------------------------------------------------------------
\begin{frame}[fragile]{Reading \file{out.dat} --- two traps}

\begin{alertblock}{Trap 1: the header lies}
The \code{\# LF =} line prints \code{Lmin}, \textbf{not} the actual final
longitude. The real $\Lf$ is $L_0 + \mathrm{par}(1)(L_{\min}-L_0)$ --- and it
is the last value in column 1 of the data block.
\end{alertblock}

\begin{alertblock}{Trap 2: one record spans twenty lines}
Each record is 20 numbers ($L$, $Y(1{:}16)$, $u(1{:}3)$), but the write format
\code{'(e24.16,2x,(e24.16e3))'} triggers Fortran \textbf{format reversion}.
Never read it line-by-line.
\end{alertblock}

\begin{lstlisting}[style=py]
vals = [float(t) for l in open('out.dat').read().split('# Data')[1].split('\n')
                 for t in l.split()]
rows = [vals[k:k+20] for k in range(0, len(vals), 20)]   # L, Y(1:16), u(1:3)
\end{lstlisting}

\code{rows[-1][1]} $=$ final $P$, \quad \code{rows[-1][7]} $=$ final mass,
\quad \code{rows[-1][8]} $= t_f$. \\
(Tested against the reference run's \file{out.dat}.)
\end{frame}

%------------------------------------------------------------------
\begin{frame}{Setting up a new case}

\begin{enumerate}
\item \textbf{Reproduce the reference run first.} If it does not match, the
      problem is your build, not your case.
\item Start from \file{in10N.dat} (cold) or a converged \file{next.dat}
      (warm --- preferred).
\item \code{par(2)} $\leftarrow$ new thrust [N].
\item \code{par(5:9)} if the target changes; \code{Y0(1:5)}, \code{Y0(7)},
      \code{L0} if the departure orbit or mass changes.
\item \alert{\code{Lmin} --- this is the one that bites.} The shipped
      $29.901174314034$ is the minimum-time final longitude \emph{for 10 N},
      and is wrong at any other thrust. Getting $L_{\min}$ at a new thrust is
      its own minimum-time problem. No shortcut inside v1.
\item \code{par(1)}: larger $=$ longer, more leisurely, higher final mass.
\item Cold start? Reset \code{lambda} to \code{0. 0.} and give \code{ZI(1)} a
      positive $t_f$ guess in hours.
\item Run, then work the four verification checks in order.
\end{enumerate}
\end{frame}

%------------------------------------------------------------------
\begin{frame}{Gotchas}

\begin{alertblock}{The one real hazard: \code{Pfun} ignores integration
failure}
rkf45 \code{iflag = 6} means \emph{``integration was not completed''}.
\code{Pfun} prints it and then uses \code{Y} as though it were at $\Lf$.
It fires \textbf{32 times} in the reference run --- which is still correct,
because all 32 sit in the $\lambda_1$ / Jacobian phase and none in the final
refinement.\\[3pt]
\textbf{Check:} where do the \code{[PFUN] IFAIL} lines sit relative to the
\code{BEFORE/AFTER REFINEMENT} blocks? Any inside or after them invalidates
the answer.
\end{alertblock}

\begin{itemize}\small
\item \code{TRACE} is a \textbf{unit number} --- \code{2} writes
      \file{fort.2}. \code{0} disables; $>1$ also turns on verbose stdout.
\item Two \textbf{silent clamps} in \code{Phifun}: $x_1 = \max(10^{-3},P)$,
      and $\mathrm{d}t = 0$ when $F(6)\le 0$. Both hide non-physical
      excursions rather than erroring.
\item \code{homCI} can bail: \code{Initial conditions homotopy failed}.
      Loosen \code{par(10)}/\code{par(11)}, or warm start.
\item \file{matlab/main.m} is \textbf{not part of the solver} --- an old
      NN-toolbox script training a five-neuron $t_f$ estimator.
\end{itemize}
\end{frame}

%------------------------------------------------------------------
\begin{frame}{v0 versus v1}

\begin{center}\small
\begin{tabular}{@{}lll@{}}
\toprule
 & v0 & v1\\
\midrule
independent variable & time $t$ & longitude $L$\\
state dimension $n$ & 7 & 8 ($t_f$ added)\\
fixed & $t_f$ \emph{and} $\Lf$ & $\Lf$ only\\
free & --- & $t_f$\\
span control & $\mathrm{par}(1)=c_{t_f}$ on $[t_0,t_{\min}]$ &
$\mathrm{par}(1)$ on $[L_0,L_{\min}]$\\
\code{lpar} / \code{lipar} & 14 / 2 & 11 / 1\\
tolerances & \code{arcre/ae} $10^{-5}$, \code{ansre/ae} $10^{-8}$ &
$10^{-6}$, $10^{-10}$ (tighter)\\
validated here & 2026-07-20: \textbf{1378.37 kg} &
2026-08-15: \textbf{1381.41 kg}, $t_f$ free\\
\bottomrule
\end{tabular}
\end{center}

\file{hybrd.f}, \file{rkf45.f}, \file{LAPACK.f} are \textbf{byte-identical}
between the two --- same solver core, different formulation.

\begin{alertblock}{Not the same problem instance}
The fixed/free roles are swapped and the span multipliers act on different
axes. Do \emph{not} read 1381.41 vs 1378.37 as a discrepancy.
\end{alertblock}
\end{frame}

%------------------------------------------------------------------
\begin{frame}{What to borrow when minimum-fuel work resumes}

\begin{enumerate}
\item \textbf{Arclength continuation for energy $\to$ fuel.} HOMPACK
      predictor--corrector with automatic steplength and re-refinement, versus
      our hand-scheduled $\varepsilon$ ladder. Candidate cure for where our
      ladder stalls --- the \file{GTO\_tulip} $1.01$--$1.11\times$
      near-minimum-time band, and the deep thrust rungs.
\item \textbf{The start-at-the-target initial-condition homotopy.} A seeding
      trick we use nowhere; it removes the need for a good cold guess.
\item \textbf{Freeing a scalar as a zero-derivative state.} Note the
      direction: v1 fixes the span and frees time, while the tulip's open
      problem needs the \emph{span} freed. The \emph{device} transfers, not
      the formulation --- carry the span as the zero-derivative state.
\item \textbf{Its converged costates as clean seeds.} \code{\# Z} in
      \file{out.dat} is PMP-converged, unlike our direct-KKT duals at high
      eccentricity. Already the recorded plan for our own MATLAB indirect
      solver.
\end{enumerate}

\begin{block}{Open cross-check}
Give v0 the same transfer with $t_f$ fixed at $146.3947$ h. If it returns
$1381.41$ kg and $\Lf = 56.66$ rad, the two formulations validate each other
\emph{and} our direct campaign. Blocker is clerical: decode v0's
$\mathrm{lpar} = 14$ \code{PAR} layout.
\end{block}
\end{frame}

\end{document}
